\subsection{Definition of Galois extensions}

THEOREM 1. --- Let N be an algebraic extension of K and \(\Gamma\) the group of K-automorphisms of N. Then the following assertions are equivalent:

\begin{enumerate}
\def\labelenumi{\alph{enumi})}
\item
  Every element of \(\mathbf{N}\) invariant under \(\Gamma\) belongs to the image of \(\mathbf{K}\) in \(\mathbf{N}\) .
\item
  N is a separable quasi-Galois extension of K.
\item
  For each \(x \in \mathbf{N}\) the minimal polynomial of \(x\) over \(\mathbf{K}\) splits in \(\mathbf{N}[\mathbf{X}]\) into a product of distinct polynomials of degree 1.
\end{enumerate}

The equivalence of b) and c) follows from the Cor. of Prop. 6 (V, p.~40) and the definition of quasi-Galois extension (V, p.~53, Def. 2). Let us identify K with its canonical image in N.

\(a) \Rightarrow c)\) : Suppose that K is the field of invariants of \(\Gamma\) . Let \(x \in \mathbf{N}\) , with minimal polynomial \(f\) over K and let A be the set of all roots of \(f\) in N. Put

\[
g (\mathbf {X}) = \prod_ {y \in \mathsf {A}} (\mathbf {X} - y).
\]

Every automorphism \(\sigma \in \Gamma\) induces a permutation of \(\mathbf{A}\) , and hence leaves invariant the coefficients of the polynomial \(g \in \mathbf{N}[\mathbf{X}]\) . We thus have \(g \in \mathbf{K}[\mathbf{X}]\) and since \(g(x) = 0\) , the polynomial \(g\) is a multiple of \(f\) in \(\mathbf{K}[\mathbf{X}]\) (V, p.~16, Th. 1). Further, \(f\) and \(g\) are monic and \(g\) divides \(f\) (IV, p.~16, Prop. 5); thus we have \(f = g\) , which means that the minimal polynomial \(f\) of \(x\) over \(\mathbf{K}\) is a product in \(\mathbf{N}[\mathbf{X}]\) of distinct polynomials of degree 1.

\(c) \Rightarrow a)\) : let \(x\) be an element of \(\mathbf{N}\) not belonging to \(\mathbf{K}\) . Denote by \(\Omega\) an algebraic closure of \(\mathbf{K}\) containing \(\mathbf{N}\) as subextension (V, p.~23, Th. 2). Let \(f\) be the minimal polynomial of \(x\) over \(\mathbf{K}\) , which is of degree \(\geqslant 2\) by hypothesis and let \(\mathbf{A}\) be the set of roots of \(f(\mathbf{X})\) in \(\mathbf{N}\) . If condition \(c)\) holds, we have \(f(\mathbf{X}) = \prod_{y \in \mathbf{A}} (\mathbf{X} - y)\) and so

(V, p.~53, Cor. 1) A is the set of conjugates of \(x\) in \(\Omega\) . Since \(f\) has degree \(\geqslant 2\) , there exists in A an element \(y \neq x\) , hence a K-automorphism \(u\) of \(\Omega\) such that \(u(x) = y\) . Now under the hypothesis \(c\) ), the extension N of K is quasi-Galois, hence \(u(N) = N\) (V, p.~54, Cor. 1); it follows that \(u\) induces a K-automorphism \(\sigma\) of N such that \(\sigma(x) = y \neq x\) , so K is the field of invariants of \(\Gamma\) .

DEFINITION 1. --- An extension N of K is said to be Galois if it is algebraic and satisfies the equivalent conditions a), b), c) of Th. 1.

Let N be a field, \(\Gamma\) a group of automorphisms of N and \(N_{0}\) the field of invariants of \(\Gamma\) . When N is algebraic over \(N_{0}\) it is a Galois extension of \(N_{0}\) . This need not always be so: for example, suppose that K is infinite and take N to be the field of rational fractions \(\mathrm{K}(\mathrm{X})\) ; for each \(a \in K\) let \(\sigma_{a}\) be the automorphism of \(\mathrm{K}(\mathrm{X})\) which maps \(f(\mathrm{X})\) to \(f(\mathrm{X} + a)\) . The set of all \(\sigma_{a}\) is a group of automorphisms of \(\mathrm{K}(\mathrm{X})\) whose field of invariants is easily seen to be K; however \(\mathrm{K}(\mathrm{X})\) is not algebraic over K.

Let \(\Omega\) be an algebraic closure of K, let A denote a set of elements of \(\Omega\) separable over K and B the set of conjugates over K of elements of A. Then B consists of elements that are algebraic and separable over K. Therefore (V, p.~39, Prop. 6 and p.~56, Prop. 5) the field K(B) is a separable quasi-Galois extension of K; in other words, the quasi-Galois extension generated by A (V, p.~55) is a Galois extension of K; we shall also say that it is the Galois extension of K generated by the subset A of \(\Omega\) .

In particular, the splitting field in \(\Omega\) of a family of separable polynomials over K, a separable closure of K, are Galois extensions of K.

PROPOSITION 1. --- Let N be an extension of K and \((\mathbf{N}_{i})_{i\in I}\) a non-empty family of subextensions of N. We put \(E=\bigcap_{i\in I}N_{i}\) and \(F=K\left(\bigcup_{i\in I}N_{i}\right)\) . If all the extensions

\(N_{i}\) are Galois over K then the same is true of E and F.

In the first place E is algebraic and separable over K (V, p.~36, Prop. 1) and the same holds for F (V, p.~41, Prop. 8). Moreover, E and F are quasi-Galois over K by Prop. 4 (V, p.~55).

PROPOSITION 2. --- Let N be a Galois extension of K and E a subextension of N, of finite degree over K. There exists a subextension F of N containing E, Galois and of finite degree over K.

Since N is quasi-Galois over K, Cor. 1 of V, p.~56 proves the existence of a quasi-Galois subextension F of N containing E and of finite degree over K. Since N is separable over K, the same holds for F (V, p.~36, Prop. 1), hence F is Galois over K.

Prop. 2 implies the following result: let \(\Omega\) be an algebraic closure of K and \(E_1, \ldots, E_n\) separable algebraic extensions of finite degree over K, contained in \(\Omega\) . Then there exists a Galois extension N of K, of finite degree, contained in \(\Omega\) and containing \(E_1, \ldots, E_n\) .

